\section{Conclusiones y Perspectivas}
El desarrollo del formalismo de la matriz denisdad  $\tilde{\rho}$ aplicado a un sistema atómico de tres niveles en configuración $\Lambda$ ha permitiddo comprender de forma estructurada y rigurosa el fenómeno de la Transparencia Inducida Electromagnéticamente (EIT) y la luz lenta. Se mostró que la interferencia cuántica destructiva entre las amplitudes de transición óptica anula de forma selectiva la componente imaginaria de la susceptibilidad eléctrica bajo la resonancia de dos fotones $\text{Im}[\chi]_{\delta = 0} \to 0$, suprimiendo la absocrión y permitiendo una ventana de transparencia en el centro de la resonancia.

Además, tanto las expresiones para los perfiles analíticos como los perfiles numéricos confirman las características espectrales de las ventanas de transparencia ocasionadas por el EIT y su dependencia direcia a la intensidad del haz de control $\Omega_c$ y la desintonía del haz de prueba $\Delta_p$. Ésto permite que el EIT se use como herramienta para regímenes de luz lenta y almacenamiento de pulsos coherentes de luz.

Si bien los modelos teóricos ideales de sistemas de tres niveles describen con precisión el fenómeno del EIT así como la forma de su perfil espectral, su implementación y validación frente a espectros experimentales reales con celdas de vapor atómico plantean diversos desafíos de gran importancia tanto académica como tecnológica.

Los ajustes de perfiles espectroscópicos convencionales (como pueden ser perfiles lorentzianos o perfiles de Voigt) resultan isnuficientes para describir la forma del perfil, elestrechamiento y el dip de transparencia que caracteriza al EIT. Por lo que la solución analítica exacta de la susceptibilidad eléctrica $\chi(\Delta_p)$ (y por tanto el coeficiente de absorción $\alpha (\Delta_p)$ y la transmisión $\mathcal{T}$) en el ajuste numérico de perfiles espectrales experimentales permite extraer con ala precisión parámetros físicos críticos de los experimentos, como pueden ser tasas de relajación, decoherencia delos estados fundamentales, corrimientos energéticos debido al efecto Stark dinámico y la calibración efectiva de freciencias de Rabi.

En este sentido la precisión en los modelados teóricos así como en los algoritmos y métodos de ajuste de las líneas espectrales resultan indispensables para el desarrollo de equipo de alta sensibilidad, referencias atómicas de frecuencia compactas y memorias cuánticas. Por lo que su investigación tanto en la teoría, el elemento computaciional como en la instrumentación reulta es de alto interés en la actualidad.
